\subsection{Principal Component Analysis}
The eigenvalue decomposition of the standardized feature covariance matrix reveals that the effective dimensionality of the data is significantly lower than 12. As shown in Table \ref{tab:pca_variance}, the first three principal components (PCs) explain 73.49\% of the total variance, and the first four explain 81.78\%. This justifies our choice of $k=4$ for the reduced dimensionality experiments, as we retain over 80\% of the signal while reducing the feature count by two-thirds.

\begin{table}[htbp]
\caption{PCA Eigenvalues and Explained Variance}
\begin{center}
\begin{tabular}{lccc}
\hline
\textbf{Component} & \textbf{Eigenvalue} & \textbf{Explained \%} & \textbf{Cumulative \%} \\
\hline
PC1 & 3.737 & 31.14 & 31.14 \\
PC2 & 2.971 & 24.76 & 55.90 \\
PC3 & 2.111 & 17.59 & 73.49 \\
PC4 & 0.994 & 8.29 & 81.78 \\
PC5 & 0.927 & 7.72 & 89.50 \\
\hline
\end{tabular}
\end{center}
\label{tab:pca_variance}
\end{table}

The Scree Plot (Figure \ref{fig:scree_plot}) and Pareto Plot (Figure \ref{fig:pareto_plot}) confirm this trend, with a distinct "elbow" appearing after the third component.

\begin{figure}[htbp]
\centering
\includegraphics[width=0.45\textwidth]{figures/ScreePlot.png}
\caption{Scree Plot showing eigenvalues for each principal component.}
\label{fig:scree_plot}
\end{figure}

\begin{figure}[htbp]
\centering
\includegraphics[width=0.45\textwidth]{figures/ParetoPlot.png}
\caption{Pareto plot of explained variance. The dashed line marks the 80\% threshold, achieved by 4 PCs.}
\label{fig:pareto_plot}
\end{figure}

Analysis of the factor loadings (Figure \ref{fig:biplot}) allows us to interpret the physical meaning of these components:
\begin{itemize}
    \item \textbf{PC1 (31.1\%):} Dominated by stellar environment variables (\texttt{koi\_teq}, \texttt{koi\_insol}, \texttt{koi\_srad}).
    \item \textbf{PC2 (24.8\%):} Heavily weighted by signal strength features (\texttt{koi\_depth}, \texttt{koi\_model\_snr}) and planetary radius (\texttt{koi\_prad}).
    \item \textbf{PC3 (17.6\%):} A mix of orbital characteristics and stellar properties (\texttt{koi\_period}, \texttt{koi\_slogg}).
\end{itemize}

\begin{figure}[htbp]
\centering
\includegraphics[width=0.45\textwidth]{figures/PCA_Biplot.png}
\caption{PCA Biplot showing the first two principal components. Vectors indicate the direction and strength of feature contributions.}
\label{fig:biplot}
\end{figure}

\subsection{Model Benchmarking}
We first benchmarked multiple classifiers using 10-fold cross-validation. Table \ref{tab:model_comparison} summarizes the results on the original 12-dimensional feature set. Among the candidate models benchmarked using PyCaret—Logistic Regression, Random Forest, MLP, and gradient boosting methods such as LightGBM \cite{lightgbm}—LightGBM and Random Forest achieved the highest cross-validated accuracy.

\begin{table}[htbp]
\caption{Top Model Performance (Original Features)}
\begin{center}
\begin{tabular}{lccccc}
\hline
\textbf{Model} & \textbf{Accuracy} & \textbf{AUC} & \textbf{Recall} & \textbf{Prec.} & \textbf{F1} \\
\hline
LightGBM & 0.8841 & 0.9509 & 0.9024 & 0.8627 & 0.8819 \\
Gradient Boosting & 0.8804 & 0.9445 & 0.9049 & 0.8544 & 0.8788 \\
Random Forest & 0.8769 & 0.9440 & 0.8817 & 0.8646 & 0.8729 \\
MLP Classifier & 0.8913 & 0.9504 & 0.9027 & 0.8772 & 0.8897 \\
Logistic Reg. & 0.8137 & 0.8668 & 0.8889 & 0.7627 & 0.8207 \\
\hline
\end{tabular}
\end{center}
\label{tab:model_comparison}
\end{table}

When trained on the 4-PC subset, performance generally decreased. For example, Random Forest accuracy dropped to $\approx 80.9\%$, and MLP to $\approx 81.0\%$. This indicates that while the first 4 PCs capture most variance, the remaining 20\% contains fine-grained information useful for classification.

\subsection{Test Set Performance}
Table \ref{tab:test_performance} presents the final evaluation of our selected models on the held-out test set. The MLP classifier trained on original features achieved the highest overall performance.

\begin{table}[htbp]
\caption{Final Test Set Performance}
\begin{center}
\begin{tabular}{llcc}
\hline
\textbf{Model} & \textbf{Features} & \textbf{Accuracy} & \textbf{F1 Score} \\
\hline
Logistic Reg. & Original (12-D) & 0.7990 & 0.8069 \\
Logistic Reg. & PCA (4 PCs) & 0.7540 & 0.7669 \\
Random Forest & Original (12-D) & 0.8773 & 0.8738 \\
Random Forest & PCA (4 PCs) & 0.8011 & 0.7949 \\
\textbf{MLP} & \textbf{Original (12-D)} & \textbf{0.8926} & \textbf{0.8898} \\
MLP & PCA (4 PCs) & 0.8129 & 0.8049 \\
\hline
\end{tabular}
\end{center}
\label{tab:test_performance}
\end{table}

\subsection{Model Diagnostics: Confusion Matrices and ROC}
To better understand the error modes of our models, we examine the confusion matrices for the three selected classifiers on the original feature set (Figure \ref{fig:confusion_matrices}).
\begin{itemize}
    \item The \textbf{MLP} (bottom) shows the most balanced performance, with relatively few false positives (classifying low-quality as high-quality).
    \item \textbf{Logistic Regression} (top) struggles significantly more, with a higher rate of misclassification for both classes, reflecting its inability to model non-linear boundaries.
\end{itemize}

\begin{figure}[htbp]
\centering
\includegraphics[width=0.45\textwidth]{figures/Confusion_Matrix_Original.png}
\caption{Confusion Matrices for Logistic Regression, Random Forest, and MLP on the test set (Original Features).}
\label{fig:confusion_matrices}
\end{figure}

The ROC curves in Figure \ref{fig:roc_original} provide further insight into the trade-off between sensitivity and specificity. The MLP and Random Forest curves hug the top-left corner more closely than Logistic Regression, confirming their superior discriminative power across different decision thresholds.

\begin{figure}[htbp]
\centering
\includegraphics[width=0.45\textwidth]{figures/ROC_Curve_Original.png}
\caption{ROC Curves for the models trained on original features.}
\label{fig:roc_original}
\end{figure}

\subsection{Decision Boundaries and Feature Importance}
Figure \ref{fig:decision_boundaries} visualizes the decision boundaries of the three models in the PC1-PC2 plane. The Logistic Regression boundary is linear and fails to capture the central cluster of high-quality candidates efficiently. In contrast, the Random Forest and MLP models learn non-linear boundaries that better encircle the target class.

\begin{figure}[htbp]
\centering
\includegraphics[width=0.45\textwidth]{figures/Decision_Boundaries.png}
\caption{Decision boundaries of LR, RF, and MLP in the PC1-PC2 latent space.}
\label{fig:decision_boundaries}
\end{figure}

Feature importance analysis from the Random Forest model (Figure \ref{fig:feat_imp}) confirms that planetary radius (\texttt{koi\_prad}) is the most critical predictor, followed by transit depth and period. This aligns with physical intuition: larger planets with deeper, clearer transits are easier to confirm.

\begin{figure}[htbp]
\centering
\includegraphics[width=0.45\textwidth]{figures/RF_Feature_Importance_Original.png}
\caption{Random Forest feature importance (Original Features).}
\label{fig:feat_imp}
\end{figure}
